%% Section 6: Preparation for counting points
\subsection{The universal family and non-degeneracy}
\label{subsec:univfamily}
In this section, we fix the basic setup to prove 
Proposition~\ref{PropAlgPtFar}, described as the alternative on page~\pageref{eq:alternative}, and
our main results.

Fix an integer $g \ge 2$.
Recall from $\mathsection$\ref{subsec:hgtineqnondeg} that
$\mathbb{M}_g$ 
denotes the fine moduli space of smooth curves of genus $g$,
with level-$\ell$-structure where $\ell\ge \bestlevel$ is fixed, \textit{cf.}
 \cite[Chapter~XVI, Theorem~2.11 (or above
Proposition~2.8)]{ACG:Curve},
 \cite[(5.14)]{DM:irreducibility}, or \cite[Theorem
 1.8]{OortSteenbrink}. It is known that 
$\mathbb{M}_g$ is a regular, quasi-projective variety of dimension
$3g-3$. We regard it over $\IQbar$; it is irreducible according
to our convention introduced below Theorem~\ref{ThmHtInequality}. 
%It is irreducible by \cite[Theorem 5.15]{DM:irreducibility}.
There exists a universal curve $\mathfrak{C}_g$ over
$\mathbb{M}_g$,  it is smooth and proper over
$\mathbb{M}_g$ and its fibers are smooth curves of genus
$g$. Moreover, $\mathfrak{C}_g\rightarrow \mathbb{M}_g$ is
projective, \textit{cf.} \cite[Corollary to Theorem
1.2]{DM:irreducibility} or  \cite[Remark 2, \S 9.3]{NeronModels}.

Denote by $\mathrm{Jac}(\mathfrak{C}_g)$ the relative Jacobian of
$\mathfrak{C}_g \rightarrow \mathbb{M}_g$. It is an abelian scheme
coming with a natural principal polarization and equipped with level-$\ell$-structure, see
%see \cite[Proposition 4, \S 9.4]{NeronModels} or
\cite[Proposition 6.9]{MFK:GIT94}.


%admitting a section $\epsilon \colon \mathbb{M}_g \rightarrow \mathfrak{C}_g$ which is defined over $\IQbar$. \marginpar{\tiny{ZG: Global section may not exist. Not a problem but we need to embed $\mathfrak{C}_g$ piece-by-piece (w.r.t. the base) in $\mathfrak{A}_g$. Since only finitely many pieces, OK.}}  Then $\epsilon$ gives rise to an Abel-Jacobi embedding $\mathfrak{C}_g \rightarrow \mathrm{Jac}(\mathfrak{C}_g)$, which is fiberwise the classical Abel-Jacobi embedding. Note that $\mathrm{Jac}(\mathfrak{C}_g) \rightarrow \mathbb{M}_g$ is the a principally polarized abelian scheme. From now on by abuse of notation, we denote by $\mathfrak{C}_g$ its image under the Abel-Jacobi embedding. Hence $\mathfrak{C}_g$ is a subvariety of $\mathrm{Jac}(\mathfrak{C}_g)$.%, which \textit{a priori} depends on $\epsilon$.

Recall from $\mathsection$\ref{subsec:hgtineqnondeg} that  $\mathbb{A}_g$ denotes the fine moduli space of principally polarized
abelian varieties of dimension $g$, with level-$\ell$-structure. Moreover,
$\mathbb{A}_g$ is  regular and quasi-projective; see
\cite[Theorem~7.9 and below]{MFK:GIT94} or \cite[Theorem 1.9]{OortSteenbrink}.
We regard it as defined over $\IQbar$;
% \Pcomment{Perhaps we should add
%   that we need to choose an $\ell$-th root of unity in $\IQbar$ once
%   and for all.}
it is irreducible according to
our convention. Let $\pi \colon \mathfrak{A}_g
\rightarrow \mathbb{A}_g$ be the universal abelian variety; it is an
abelian scheme. Note that $\pi$ is projective; we refer to
Remark~\ref{rmk:projembedding} for this and other details.

As $\mathbb{A}_g$ is a fine moduli space we have the following Cartesian diagram
\begin{equation}
\label{EqUnivJac}
\begin{aligned}
\xymatrix{
\mathrm{Jac}(\mathfrak{C}_g) \ar[r] \ar[d]  \pullbackcorner & \mathfrak{A}_g \ar[d]^{\pi} \\
\mathbb{M}_g \ar[r]_{\tau} & \mathbb{A}_g
}
\end{aligned}
\end{equation}
the bottom arrow is the Torelli morphism. As we have level structure,
the Torelli morphism need not be injective on $\IC$-points, but it is
finite-to-one on such points, \textit{cf.} \cite[Lemma 1.11]{OortSteenbrink}. 
%% \Pcomment{A warning on Torelli. The
%% classical theorem implies injectivity on the coarse moduli space. I
%% think its generically $2-$to$-1$ with level structure $\ge 2$ and if
%% $g\ge 3$.} 
%Assume from now on that $\ell$ is even. 

% \Pinlinecomment{Danger: 
  % $\cL_s$ does not induce a principal polarization on the fibers
  % $\mathfrak{A}_{g,s}$. So $\cL$ does not induce the polarization
  % mentioned in \texttt{(HYP)}.}


%It is known that $[-1]^*\cL = \cL$; see \cite[Proposition~10.8 and 10.9]{PinkThesis}. 

We also fix an ample line bundle $\cM$ on $\overline{\mathbb{A}_g}$,
where  $\overline{\mathbb{A}_g}$ is a, possibly non-regular,
projective  variety containing $\mathbb{A}_g$ as a Zariski open and
dense subset. The Height Machine
provides an equivalence class of height functions of which we fix a
representative $h_{\overline{\mathbb{A}_g},\cM} \colon
\overline{\mathbb{A}_g}(\IQbar)\rightarrow \IR$.

%% [PH] Omitted the below, probably not necessary to use Torelli here.
%% In what follows, we will see $\mathrm{Jac}(\mathfrak{C}_g)$ as a subvariety of $\mathfrak{A}_g$, and will not distinguish $\mathbb{M}_g$ and its image in $\mathbb{A}_g$ under the morphism on the bottom. For example when we talk about subvarieties of $\mathbb{M}_g$, we see them simultaneously as subvarieties of $\mathbb{M}_g$ and of $\mathbb{A}_g$. Note that $\dim \mathbb{M}_g = 3g-3$ for each $g \ge 2$.


%Denote as usual by $\mathbb{M}_g$ the open Torelli locus, namely the image of $\mathbb{M}_g$ under the bottom morphism. Then $\dim \mathbb{M}_g = 3g-3$ for each $g \ge 2$.


Next we fix a projective embedding of $\mathfrak{A}_g$ over
$\mathbb{A}_g$.
There is a relatively ample line bundle
$\cL$ on $\mathfrak{A}_g/\mathbb{A}_g$ with $[-1]^*\cL = \cL$; see
\cite[Th\'{e}or\`{e}me~XI~1.4]{LNM119}. 
After replacing $\cL$
by $\cL^{\otimes N}$, with $N\ge 4$ large enough, we can assume that
$\cL$ is very ample relative over $\mathbb{A}_g$ and $[-1]^*\cL=\cL$.
By \cite[Proposition~4.4.10(ii) and Proposition~4.1.4]{EGAII}, we
then have a closed immersion $\mathfrak{A}_g \rightarrow
\IP^n_{\IQbar} \times \mathbb{A}_g$ over $\mathbb{A}_g$
arising from $\cL \otimes
\pi^*(\cM|_{\mathbb{A}_g}^{\otimes p})$
for some integer $p\ge 1$, note that $\cM|_{\mathbb{A}_g}$ is ample.
%for some very ample line bundle $\cM$ on $\mathbb{A}_g$. 
%In this way we obtain an immersion for each fiber of $\mathfrak{A}_g$ and a global immersion into some projective bundle. We would rather have a global immersion into a product of $\IP^n_\IQbar$ with a variety. There is a finite collection of  Zariski open affine subsets $V_1,\ldots,V_t$ of $\mathbb{A}_g$ such that $\pi_*\cL|_{V_i}$ is a free sheaf of modules for all $1\le i\le t$. So $\mathbb{P}(\pi_*\cL|_{V_i})\cong \IP_\IQbar^n\times V_i$ by flat base change. Moreover, for each $i$ we get, by restriction,  a closed immersion $\mathfrak{A}_{V_i}=\pi^{-1}(V_i)\colon \mathfrak{A}_{V_i}\rightarrow \IP^n_\IQbar \times V_i$.
For each $s \in \mathbb{A}_g(\IQbar)$, the fiber
$\mathfrak{A}_{g,s}=\pi^{-1}(s)$ is realized as a
projective subvariety of $\IP^n_\IQbar$ and the induced
closed immersion $\mathfrak{A}_{g,s}\rightarrow\IP^n_\IQbar$ comes
from the restriction $\cL|_{\mathfrak{A}_{g,s}}$ which is ample.
Flatness of $\mathfrak{A}_g\rightarrow\mathbb{A}_g$
implies that
$\dim H^0(\mathfrak{A}_{g,s},\cL_{\mathfrak{A}_{g.s}})$ is independent of $s$.
So $\mathfrak{A}_{g,s}$ is projectively
normal inside $\IP^n_\IQbar$, a property that will play a role later on. 
%\footnote{
%  Here is an explanation why we get a fiberwise projectively normal
%  immersion. To be removed later.
%  Let $s\in \mathbb{M}_g$.
%  We abuse our convention and write $\mathfrak{A}_{g,s}$
%  for $\mathfrak{A}_g\times_S \mathrm{Spec\,}k(s)$ and consider it as
%  an abelian variety defined over $k(s)$. 
%  Then $\cL$ restricts to an ample line bundle $\cL_s$ on
%  $\mathfrak{A}_{g,s}$. As $\cL_s$ is the $N$-th power of an ample
%  line bundle on $\mathfrak{A}_{g,s}$ and as $N\ge 4$ we have
%  $H^0(\mathfrak{A}_{g,s},\cL_s)\not=0$. 
%  So $H^p(\mathfrak{A}_{g,s},\cL_s)=0$ for all $p\ge 1$
%  by~\cite[Vanishing Theorem, Section III.16]{MumfordAbVar}.
%  By~\cite[Corollary, Section II.5]{MumfordAbVar}
%  the Euler characteristic $\sum_{p\ge 0}(-1)^p \dim_{k(s)}
%  H^p(\mathfrak{A}_{g,s},\cL_s)=H^0(\mathfrak{A}_{g,s},\cL_s)$ is locally constant on
% $\mathbb{M}_g$; here we need that
%  $\mathfrak{A}_g\rightarrow\mathbb{M}_g$ is proper and
%  flat. So $s\mapsto
%  H^0(\mathfrak{A}_{g,s},\cL_s)$ is constant with value $n+1$, say, on
%  the connected scheme
%  $\mathbb{M}_g$.
%    Then
%  $\pi_*\cL$ is a locally free sheaf of finite type  on $\mathbb{A}_g$
% with local rank $n+1$ and    $s^*\pi_*\cL$ is naturally isomorphic to
%  $H^0(\mathfrak{A}_{g,s},\cL_s)$ for all $s$ by \cite[Corollary and
%  Lemma 1~II.5.2]{MumfordAbVar}.
% Consider the immersion
%  $\mathfrak{A}_g \rightarrow\mathbb{P}(\pi_*\cL)$ over $\mathbb{A}_g$
%  induced by the natural surjective  homomorphism
%  $\pi^*\pi_*\cL\rightarrow \cL$, \textit{cf.} \cite[Proposition~13.56]{GoertzWedhorn}. It
%  is a closed immersion as $\pi$ is proper. On each fiber we obtain a
%  closed immersion $\mathfrak{A}_{g,s}\rightarrow
%  \mathbb{P}(s^*\pi_*\cL) \cong
%  \mathbb{P}(H^0(\mathfrak{A}_{g,s},\cL_s))$; here the bundle $\cO(1)$
%  on the projective space pulls back to $\cL_s$. As $N\ge 4$ the
%  abelian variety $\mathfrak{A}_{g,s}$ is projectively normal in
%  $\mathbb{P}(H^0(\mathfrak{A}_{g,s},\cL_s))$.}
% For each $s \in
% \mathbb{A}_g(\IQbar)$, the line bundle $\cL$ restricted to the abelian
% variety $\mathfrak{A}_{g,s}$ is ample. 
% \Pcomment{Think about
% projectively normal.}
% \Pcomment{I uncommented this sentence to define the notation
  % $\mathfrak{A}_{g,s}$ used later on. Maybe we can drop the statement
  % on ampleness.}



Recall that $\cL$ is symmetric and very ample on each fiber of
$\mathfrak{A}_{g}$. 
By Tate's Limit Argument we obtain the fiberwise N\'eron--Tate height,
\textit{cf.} (\ref{eq:fiberwiseNT}), 
% $\mathfrak{A}_g(\IQbar)\rightarrow [0,\infty)$. These height
%   functions match on the overlap and we obtain the fiberwise
%   N\'eron--Tate height
\begin{equation}
 \label{eq:NTfinemoduli}
  \hat h \colon \mathfrak{A}_{g}(\IQbar)\rightarrow [0,\infty).
\end{equation}

%% Instead of twisting by a power of $\pi^*(\cM|_{\mathbb{A}_g})$ we
%% could have instead embedded $\mathfrak{A}_g\rightarrow \mathbb{A}_g$
%% over members of a finite open affine cover of $\mathbb{A}_g$. This
%% would leads to the same global N\'eron--Tate height 
%% as $\pi^*(\cM|_{\mathbb{A}_g})$ restricted to a single
%% $\mathfrak{A}_{g,s}$ is trivial. 



Let $M\ge 1$ be an integer. We write $\mathfrak{A}_g^{[M]}$ for the
$M$-fold fibered power
$\mathfrak{A}_g\times_{\mathbb{A}_g} \cdots \times_{\mathbb{A}_g}\mathfrak{A}_g$
over $\mathbb{A}_g$.
Then $\mathfrak{A}_g^{[M]}\rightarrow \mathbb{A}_g$ is an abelian
scheme.
% Similarly, we define $\mathfrak{C}_g^{[M]}$ to be the $M$-fold fibered
% power $\mathfrak{C}_g\times_{\mathbb{M}_g}\cdots
% \times_{\mathbb{M}_g}\mathfrak{C}_g$.


By taking the product we obtain closed
immersions $\mathfrak{A}_g^{[M]}\rightarrow (\IP_\IQbar^n)^{M}\times
\mathbb{A}_g$. The fiber of $\mathfrak{A}_g^{[M]}\rightarrow
\mathbb{A}_g$ above $s\in \mathbb{A}_g(\IQbar)$ is the $M$-fold power of
$\mathfrak{A}_{g,s}$. The associated fiberwise N\'eron--Tate height
$\hat h \colon \mathfrak{A}_g^{[M]}(\IQbar)\rightarrow [0,\infty)$ is
the sum of the N\'eron--Tate heights, as in (\ref{eq:NTfinemoduli}),
of the $M$ coordinates.%\Pcomment{Replaced $m$ by $M$ twice (typo).}

Let us now define the Faltings--Zhang morphism. 
In our setting the relative Picard scheme
$\mathrm{Pic}(\mathfrak{C}_g/\mathbb{M}_g)$ exists as a group scheme
over $\mathbb{M}_g$. It is a union over all $p\in\IZ$
of open and closed subschemes
$\mathrm{Pic}^p(\mathfrak{C}_g/\mathbb{M}_g)$, where $p$ indicates the
degree of a line bundle.
By definition we have
$\mathrm{Jac}(\mathfrak{C}_g) =
\mathrm{Pic}^0(\mathfrak{C}_g/\mathbb{M}_g)$.
We cannot expect to
have a section of $\mathfrak{C}_g\rightarrow\mathbb{M}_g$,
so we cannot expect to find an immersion of $\mathfrak{C}_g$
into $\mathrm{Jac}(\mathfrak{C}_g/\mathbb{M}_g)$. 
As constructed in the proof of
\cite[Proposition~6.9]{MFK:GIT94} we
do have a morphism $\mathfrak{C}_g \rightarrow
\mathrm{Pic}^1(\mathfrak{C}_g/\mathbb{M}_g)$ over $\mathbb{M}_g$.
Let $\mathfrak{C}_g^{[M]}$ and
$\mathrm{Pic}^p(\mathfrak{C}_g/\mathbb{M}_g)^{[M]}$ denote the
respective $M$-th fibered powers over $\mathbb{M}_g$. 
The difference morphism coming from the group scheme law
$\mathrm{Pic}(\mathfrak{C}_g/\mathbb{M}_g)\times_{\mathbb{M}_g}
\mathrm{Pic}(\mathfrak{C}_g/\mathbb{M}_g)\rightarrow
\mathrm{Pic}(\mathfrak{C}_g/\mathbb{M}_g)$ restricts to a morphism
$\mathrm{Pic}^1(\mathfrak{C}_g/\mathbb{M}_g)\times_{\mathbb{M}_g}
\mathrm{Pic}^1(\mathfrak{C}_g/\mathbb{M}_g)
\rightarrow
\mathrm{Jac}(\mathfrak{C}_g/\mathbb{M}_g)$ of schemes
over $\mathbb{M}_g$.
We take the appropriate  product morphism over $\mathbb{M}_g$ to get a
morphism 
\begin{equation}
  \label{eq:prefaltingszhang}
  \mathfrak{C}_g^{[M+1]} \rightarrow
 \mathrm{Jac}(\mathfrak{C}_g/\mathbb{M}_g)^{[M]}
\end{equation}
over $\mathbb{M}_g$. The choice of product is modeled after
(\ref{eq:faltingszhangintro}). More precisely, consider the situation
above a $k$-point of $\mathbb{M}_g$, where $k$ is an algebraically
closed field. The fiber of $\mathfrak{C}_g \rightarrow \mathbb{M}_g$ above this point is a smooth curve $C$
defined over $k$ of genus $g$. For $P_0,\ldots,P_M\in C(k)$ the
morphism (\ref{eq:prefaltingszhang}) maps $(P_0,P_1,\ldots,P_M)
\mapsto (P_1-P_0,P_2-P_0,\ldots,P_M-P_0)$ where the difference takes
place in the Jacobian of $C$. %\Gcomment{Slightly rephrased this.}

Recall (\ref{EqUnivJac}). We take the $M$-fold product and compose
with (\ref{eq:prefaltingszhang}) to obtain a commutative diagram of
morphisms of schemes %\Ginlinecomment{Removed the notation $\mathscr{D}_M$ in the diagram.}
\begin{equation}
  \label{DefFaltingsZhang}
  \begin{aligned}
    \xymatrix{
      \mathfrak{C}_g^{[M+1]}\ar[r] \ar[d] &
      \mathfrak{A}_g^{[M]}
      \ar[d] \\  
      \mathbb{M}_g \ar[r]_{\tau} & \mathbb{A}_g.
    }
  \end{aligned}
\end{equation}

If $S\rightarrow \mathbb{M}_g$ is a morphism of schemes then we define
$\mathfrak{C}_S = \mathfrak{C}_g \times_{\mathbb{M}_g} S$ and
$\mathfrak{C}_S^{[M]} = \mathfrak{C}_g^{[M]}\times_{\mathbb{M}_g} S$.
If $S$ is irreducible, then so is $\mathfrak{C}_S^{[M]}$ by induction
on $M$ and a topological argument using that
$\mathfrak{C}_S\rightarrow S$ is smooth and hence open.
Taking the fibered product with $S$  and composing with (\ref{DefFaltingsZhang})
yields a commutative diagram of morphisms of schemes
\begin{equation*}
%  \label{DefFaltingsZhang2}
  \begin{aligned}
    \xymatrix{
      \mathfrak{C}_S^{[M+1]}\ar[r] \ar[d] &
      \mathfrak{A}_g^{[M]}
      \ar[d] \\  
      S\ar[r]_{\tau\circ(S\rightarrow\mathbb{M}_g)} & \mathbb{A}_g.
    }
  \end{aligned}
\end{equation*}
% we call it Faltings--Zhang morphism and also denote it by
% $\mathscr{D}_M$.
% \Pcomment{Here's a second variant of Faltings--Zhang defined on a
%   different source.}
% There is no reason to expect $\mathscr{D}_M$ to be proper. Here is an
% alternative construction that is proper.
% We continue with
% (\ref{DefFaltingsZhang2}) and use the
By the
universal property
of the fibered product we get a morphism of schemes 
\begin{equation}
  \label{DefFaltingsZhang3}
  \begin{aligned}
    \xymatrix{
      \mathfrak{C}_S^{[M+1]}\ar[r]^-{\mathscr{D}_{M}}
\ar[dr]  &
      \mathfrak{A}_g^{[M]}      \times_{\mathbb{A}_g} S
      \ar[d] \\  
      & S
    }
  \end{aligned}
\end{equation}
over $S$. We call $\mathscr{D}_{M}$ the \textit{Faltings--Zhang morphism} (over
$S$).
Then $\mathscr{D}_{M}$  is proper since the diagonal arrow in
(\ref{DefFaltingsZhang3}) is proper.

% Finally, let us consider
% the morphism
% $$\mathfrak{A}_g^{[M+1]} \rightarrow \mathfrak{A}_g^{[M]}$$
% defined fiberwise over $\mathbb{A}_g$ by
% \begin{equation}
% \label{eq:faltingszhang0}
% (P_0,P_1,\ldots,P_M) \mapsto (P_1-P_0,\ldots,P_M-P_0).
% \end{equation}
%  %This morphism is over $\IA_g$ and it is proper. 
% In fact, we will need it as a morphism on $\mathfrak{C}_g^{[M+1]}$. 
% Note that  there is no canonical embedding
% of $\mathfrak{C}_g$ in $\mathfrak{A}_g$ as we may have no section to
% construct the Abel--Jacobi map. 
% However,
% we do have a morphism
% $\mathfrak{C}_g \rightarrow
% \mathrm{Pic}^1(\mathfrak{C}_g/\mathbb{M}_g)$ to the line bundles of
% degree $1$;
% see the proof of \cite[Proposition~6.9]{MFK:GIT94}. 
% The difference of two sections of
% $\mathrm{Pic}^1(\mathfrak{C}_g/\mathbb{M}_g)$
% lies in $\mathrm{Pic}^0(\mathfrak{C}_g/\mathbb{M}_g) = \mathrm{Jac}(\mathfrak{C}_g)$.
% For each integer $M\ge 2$  the map
% (\ref{eq:faltingszhang0}) yields a well-defined morphism
% \begin{equation}
% \label{DefFaltingsZhang}
% \mathscr{D}_M\colon \mathfrak{C}_g^{[M+1]}\rightarrow \mathfrak{A}_g^{[M]}
% \end{equation}
% called the  \textit{$M$-th Faltings--Zhang} morphism.
% By abuse of notation we also write $\mathscr{D}_M$ for the 
% morphism $\mathfrak{C}_S^{[M+1]}\rightarrow \mathfrak{A}_g^{[M]}
% \times_{\mathbb{A}_g} S$
% induced by the base change by
% $S\rightarrow \mathbb{M}_g\xrightarrow{\tau} \mathbb{A}_g$.

Let for the moment $S\rightarrow \mathbb{M}_g$ be the identity.
If $s \in \mathbb{M}_g(\IQbar)$, then $\mathfrak{C}_s$ is the curve
parametrized by $s$, and $\mathfrak{A}_{g,\tau(s)}$ is its Jacobian.
To embed $\mathfrak{C}_s$ into $\mathfrak{A}_{g,\tau(s)}$ we must work with a
base point $P\in
\mathfrak{C}_s(\IQbar)$. Then 
$\mathfrak{C}_s-P = \mathscr{D}_1(\{P\} \times \mathfrak{C}_s)$
is an irreducible curve inside $\mathfrak{A}_g$ lying above $\tau(s)$.
Hence it provides a closed immersion   $\mathfrak{C}_s-P \subset
\IP_\IQbar^n$.%\Pcomment{We have to adapt this paragraph if we work
%  with the second FZ morphism.}\Gcomment{I think that this paragraph is OK because $\mathfrak{A}_g^{[M]} \times_{\mathbb{A}_g}\{s\}$ is just one abelian variety, and hence equals the fiber $\mathfrak{A}_{g,\tau(s)}^{[M]}$.}

%Say $i\in \{1,\ldots,t\}$ such that $\tau(s) \in V_i(\IQbar)$. Then the construction above provides a closed immersion   $\mathfrak{C}_s-P \subset \IP_\IQbar^n$.
% If we are given $i$ such that $s\in S_i(\IQbar)$, then
% we may consider $\mathfrak{C}_s-Q$ as an irreducible curve in $\IP^n_\IQbar$.

Let $\deg X$ denote the degree of an irreducible closed subvariety $X$ of
$\IP^n_\IQbar$ and let $h(X)$ denote its height, \textit{cf.} \cite{BGS}.

%\Pinlinecomment{TODO: Get degree bound for $P_s$ replaced by any $P$.}
\begin{lemma}
  \label{lem:uniformdegreebound}
  There exists a constant $c$ such that the following two properties
  hold for all $s\in \mathbb{M}_g(\IQbar)$.
  \begin{enumerate}
  \item [(i)] We have
    $\deg(\mathfrak{C}_s-P) \le c$
    for all $P\in
    \mathfrak{A}_{g,\tau(s)}(\IQbar)$.
  \item[(ii)]
    There exists $P_s\in \mathfrak{C}_s(\IQbar)$ such that
    $h(\mathfrak{C}_s-P_s) \le  c \max\{1,
    h_{\overline{\mathbb{A}_g},\cM}(\tau(s))\}$.    
  \end{enumerate}
  % such that if  $s\in \mathbb{M}_g(\IQbar)$, then 
  % there exists $P_s\in \mathfrak{C}_s(\IQbar)$
  % such that\Pcomment{TODO: we should now argument with $s\in
  %   S(\IQbar)$ and not $s\in \mathbb{M}_g(\IQbar)$.} \Gcomment{Maybe we only need to work with $s \in \mathbb{M}_g(\IQbar)$ as from Section 7 $S$ is always contained in $\mathbb{M}_g$. However to ``prove'' Thm 6.2, we need to first of all take a finite covering $S'$ of $S$, and thus it's better to state Thm 6.2 under the setup $S\rightarrow \mathbb{M}_g$.}
  % \begin{equation*}
  %   \deg(\mathfrak{C}_s-P_s)\le c\quad\text{and}\quad 
  %   h(\mathfrak{C}_s-P_s) \le  c \max\{1,
  %   h_{\overline{\mathbb{A}_g},\cM}(\tau(s))\}.
  % \end{equation*}
\end{lemma}
\begin{proof}
%  !!!TODO: make sure this lemma now works in the new notation.
  We need a quasi-section of
  $\mathfrak{C}_g\rightarrow\mathbb{M}_g$ as provided by \cite[Corollaire~17.16.3(ii)]{EGAIV}. 
  So there is an affine scheme $S$
  and a morphism $S\rightarrow \mathfrak{C}_g$ that factors through a
  surjective, quasi-finite, \'etale morphism $ S\rightarrow
  \mathbb{M}_g$. We consider the product
  $\mathfrak{C}_g\times_{\mathbb{M}_g} S\rightarrow
  \mathfrak{C}_g^{[2]}$ composed with the Faltings--Zhang morphism
  $\mathscr{D}_1\colon
  \mathfrak{C}_g^{[2]}\rightarrow\mathfrak{A}_g\times_{\IA_g}
  \mathbb{M}_g$ over $\mathbb{M}_g$ and then the projection of
  $\mathfrak{A}_g$. This is a morphism of schemes
  $\mathfrak{C}_g\times_{\mathbb{M}_g} S\rightarrow
  \mathfrak{A}_g$. Its 
  image is a constructible subset of $\mathfrak{A}_g$. 
  So it is a union of finitely many 
  irreducible Zariski locally closed subsets $\{X_i\}_i$ of
  $\mathfrak{A}_g$. We have
  the following property. %\Pinlinecomment{I slightly changed the
 %   argument here.}

  Given a point $s\in \mathbb{M}_g(\IQbar)$, there is an $i$ such that
  the fiber of $\pi|_{X_i}\colon X_i \rightarrow\mathbb{A}_g$ above
  $\tau(s)$ is a finite union of irreducible curves, up-to finitely
  many points one of these curves is
  $\mathfrak{C}_s-P_s$ with $P_s\in \mathfrak{C}_s(\IQbar)$.
  
  
  We  have a closed immersion $\mathfrak{A}_g\rightarrow \IP_\IQbar^n
  \times \mathbb{A}_g$.
  Moreover, a sufficiently  large positive power of $\cM$ induces a
  closed immersion of $\overline{\mathbb{A}_g}\rightarrow\IP_\IQbar^m$
   for some $m$.
  Thus, we consider $\mathbb{A}_g$ as a 
  Zariski locally closed subset
  $\IP_\IQbar^m$. We identify each 
  $X_i$ with its image in $\IP^n_\IQbar\times
  \IP^m_\IQbar$, an irreducible  Zariski locally closed set.  Then
  $\mathfrak{C}_s-P_s \subset\IP^n_\IQbar$ arises as an irreducible
  component of the intersection of some Zariski
  closure $\overline{X_i}$
   with $\IP^n_\IQbar\times
  \{\tau(s)\}$.

  We use the Segre embedding
  $\IP_\IQbar^n\times\IP_\IQbar^m\rightarrow \IP_\IQbar^{(n+1)(m+1)-1}$
  to embed our situation into projective space. By B\'ezout's Theorem
  \cite[Example~8.4.6]{Fulton},
  $\deg(\mathfrak{C}_s-P_s)$ is bounded from above uniformly in $s$.
  Translating a curve inside $\mathfrak{A}_{g,\tau(s)}$ by a point of
  $\mathfrak{A}_{g,\tau(s)}(\IQbar)$ does not
  change its degree. So if $P\in \mathfrak{A}_{g,\tau(s)}$, then
  $\deg(\mathfrak{C}_s-P)=\deg(\mathfrak{C}_s-P_s)$.  
  This yields (i).

  Part (ii) follows as (i) but this time we use
  the Arithmetic B\'ezout Theorem, still executing the intersection after
  applying the Segre embedding. Indeed, recall that
  $\mathfrak{C}_s-P_s$ as an irreducible
  of the intersection of some $\overline{X_i}$ with
  $\IP_\IQbar^n\times\{\tau(s)\}$. 
  The height and degree of $\overline{X_i}$ are
  bounded from above independently of $s$; the same holds for the
  degree of $\IP_\IQbar^n \times\{\tau(s)\}$. The height of
  $\IP_\IQbar^n \times\{\tau(s)\}$ is bounded from above linearly in terms of
  $h(\tau(s))$. Finally, we can apply \cite[Th\'eor\`eme
  3]{HauteursAlt3}.
  Finally, note that  by the Height
  Machine the absolute logarithmic Weil height $h( \tau(s))$,
  where $\tau(s)$ is understood as an element of
  $\IP^m_{\IQbar}(\IQbar)$, is bounded from
  above linearly in terms of $h_{\overline{\mathbb{A}_g},\cM}(\tau(s))$. 
\end{proof}

\subsection{Non-degeneracy of $\mathscr{D}_M(\mathfrak{C}_S^{[M+1]})$
  for large $M$}  
  In this subsection all varieties  are defined over the field $\IC$. 
We keep the notation of the previous subsection and let
$S$ be an irreducible variety with a quasi-finite %\Pcomment{Changed
%  ``finite'' to ``quasi-finite''. Is this OK?}
   morphism
$S\rightarrow\mathbb{M}_g$. %In this subsection, as in \cite{GaoBettiRank} we take $\IC$ as the base field, \textit{i.e.}, all varieties and morphisms are (viewed as) defined over $\IC$. 
Note that
$\mathscr{D}_M(\mathfrak{C}_S^{[M+1]})$ is Zariski closed
in $\mathfrak{A}_g^{[M]}\times_{\mathbb{A}_g} S$ 
because $\mathscr{D}_M$ is proper. We endow this image with the
reduced induced scheme structure.

The following non-degeneracy theorem proved by the second-named author
is crucial to prove our main result.
It confirms that
Theorem~\ref{ThmHtInequality} can be applied to
$\mathscr{D}_M(\mathfrak{C}_S^{[M+1]})$ for $M \ge 3g-2$.
We obtain a height inequality on a Zariski open dense subset. 
%\Pcomment{I added the comment on the reduced induced structure, as
%  proposed by referee .4. FZ is proper as we're using the second
%  definition.}

\begin{theorem}[\!\!{\cite[Theorem~1.2']{GaoBettiRank}}]\label{ThmNonDegForBd}
Let $S$ be an irreducible variety with a (not necessarily dominant) quasi-finite morphism $S \rightarrow \mathbb{M}_g$. Assume $g \ge 2$ and $M
\ge 3g-2$. Then $\mathscr{D}_M( \mathfrak{C}_S^{[M+1]} )$, which is a closed 
irreducible subvariety of $\mathfrak{A}_g^{[M]}\times_{\mathbb{A}_g} S$, is 
non-degenerate in the sense of
Definition~\ref{DefinitionNonDegenerate}.
\end{theorem}

The fibered product in the theorem involves
$S\rightarrow\mathbb{M}_g\xrightarrow{\tau} \mathbb{A}_g$.

More precisely, the meaning of the conclusion of the theorem is as follows. For the abelian scheme $\pi \colon \cA = \mathfrak{A}_g^{[M]} \times_{\mathbb{A}_g}S \rightarrow S$ and for the irreducible subvariety $X:=\mathscr{D}_M( \mathfrak{C}_S^{[M+1]} )$ of $\cA$, there exists a open non-empty subset $\Delta$ of $\an{S}$, with the Betti map $b_{\Delta} \colon \cA_{\Delta} = \pi^{-1}(\Delta) \rightarrow \mathbb{T}^{2g}$, such that 
\begin{equation}\label{EqNonDegXDMC}
\mathrm{rank}_{\IR} (\mathrm{d}b_{\Delta}|_{\sman{X}})_x = 2\dim X\quad \text{for some }x \in \sman{X} \cap \cA_{\Delta},
\end{equation}
when $g \ge 2$ and $M \ge 3g-2$.


\begin{proof}
This theorem, essentially \cite[Theorem~1.2']{GaoBettiRank}, is a consequence of Theorem~1.3 of \textit{loc.cit.} Because of its importance to the current paper, we hereby give more details of the deduction.

We start by showing the result for the case where $\mathfrak{C}_S
\rightarrow S$ admits a section $\epsilon$. %\Pcomment{I think we show  something stronger here, namely that 6.2 holds under the weake  assumption that $S\rightarrow\mathbb{M}_g$ is quasi-finite (but with  the additional assumption on the section). Is this correct? And, if yes, does Thm 6.2 hold with $S\rightarrow \mathbb{M}_g$ quasi-finite?, If so, we should state this.}
In this case $\epsilon$ induces an Abel--Jacobi embedding $j_{\epsilon} \colon \mathfrak{C}_S \rightarrow \mathrm{Jac}(\mathfrak{C}_S/S)$, which is a closed immersion of $S$-schemes. The modular map is the Cartesian diagram
\[
\xymatrix{
\mathrm{Jac}(\mathfrak{C}_S/S) \ar[r]^-{\iota} \ar[d] \pullbackcorner & \mathfrak{A}_g \ar[d] \\
S \ar[r] & \mathbb{A}_g
}
\]
with the bottom morphism being the composite of the given $S \rightarrow \mathbb{M}_g$ with the Torelli map $\tau \colon \mathbb{M}_g \rightarrow \mathbb{A}_g$. The Torelli map $\tau$ is quasi-finite; see \cite[Lemma 1.11]{OortSteenbrink}. Thus the bottom morphism is quasi-finite. Hence $\iota$ is quasi-finite.

We wish to apply \cite[Theorem~1.3]{GaoBettiRank} to the subvariety $j_{\epsilon}(\mathfrak{C}_S)$ of the abelian scheme $\mathrm{Jac}(\mathfrak{C}_S/S) \rightarrow S$. We need to verify the hypotheses. First of all $\iota|_{j_{\epsilon}(\mathfrak{C}_S)}$ is generically finite because $\iota$ is quasi-finite. Hypothesis (a) is satisfied since $\dim j_{\epsilon}(\mathfrak{C}_S) = \dim S + 1 > \dim S$. For hypotheses (b) and (c), note that for any $s\in S(\IC)$, the fiber $j_{\epsilon}(\mathfrak{C}_S)_s$ is the Abel--Jacobi embedding of $\mathfrak{C}_s$ in its Jacobian via the point $\epsilon(s)$. Thus hypothesis (b) is satisfied because each curve generates its Jacobian, and hypothesis (c) holds true since $g \ge 2$.

Thus we can apply \cite[Theorem~1.3.(ii)]{GaoBettiRank} and obtain that $\mathscr{D}_M(\mathfrak{C}_S^{[M+1]})$ is non-degenerate\footnote{Observe that $\mathscr{D}_M(\mathfrak{C}_S^{[M+1]}) = \mathscr{D}_M^{\cA}(j_{\epsilon}(\mathfrak{C}_S^{[M+1]}))$, with $\mathscr{D}_M^{\cA}$ be as in \cite[Theorem~1.3.(ii)]{GaoBettiRank} with $\cA = \mathfrak{A}_g\times_{\mathbb{A}_g}S \cong \mathrm{Jac}(\mathfrak{C}_S/S)$. See below \eqref{eq:prefaltingszhang}.} if $M \ge j_{\epsilon}(\mathfrak{C}_S) = \dim S + 1$. But $\dim S \le \dim \mathbb{M}_g = 3g-3$. Hence  $\mathscr{D}_M(\mathfrak{C}_S^{[M+1]})$ is non-degenerate if $M \ge 3g-2$.

%As $M \ge 3g-2 = \dim \mathbb{M}_g +1 \ge \dim S+1$, \cite[Theorem~1.2']{GaoBettiRank} implies that the theorem holds true if $\mathfrak{C}_S \rightarrow S$ admits a section.

For an arbitrary $S$, the generic fiber of $\mathfrak{C}_S \rightarrow
S$ has a rational point over some finite extension of $K(S)$, the
function field of $S$. Thus there exists a quasi-finite \'{e}tale
dominant (not necessarily surjective) morphism $\rho \colon S'
\rightarrow S$, with $S'$ irreducible, such that $\mathfrak{C}_{S'} = \mathfrak{C}_S \times_S
S' \rightarrow S'$ admits a section. Thus $X' :=
\mathscr{D}_M(\mathfrak{C}_{S'}^{[M+1]})$, as a subvariety of
$\cA':=\mathfrak{A}_g^{[M]}\times_{\mathbb{A}_g} S'$,  is
non-degenerate by the previous case. So there exists a connected, open non-empty subset $\Delta'$ of $\an{S'}$, with the Betti map $b_{\Delta'} \colon \cA'_{\Delta'} \rightarrow \mathbb{T}^{2g}$, such that for some $x' \in \sman{X'} \cap \cA'_{\Delta'}$ we have
\[
\mathrm{rank}_{\IR} (\mathrm{d}b_{\Delta'}|_{\sman{X'}})_{x'} = 2\dim X'.
\]
We may furthermore shrink $\Delta'$ so that $\rho|_{\Delta'}$ is a diffeomorphism. In particular $\Delta:=\rho(\Delta')$ is open in $\an{S}$.

Denote by $\rho'_{\cA} \colon \cA' = \cA\times_S S' \rightarrow \cA$
the projection to the first factor. Then
$\rho'_{\cA}|_{\cA'_{\Delta'}}$ is a diffeomorphism. Both $\cA
\rightarrow S$ and $\cA' \rightarrow S'$ carry
level-$\ell$-structures. By construction and uniqueness properties of
the Betti map, we may assume that $b_{\Delta} \colon \cA_{\Delta} \rightarrow \mathbb{T}^{2g}$ equals $b_{\Delta'} \circ  (\rho'_{\cA}|_{\cA'_{\Delta'}})^{-1}$. Thus
\[
\mathrm{rank}_{\IR} (\mathrm{d}b_{\Delta}|_{\sman{X}})_{x} = 2\dim X'
\]
with $x = \rho_{\cA}(x')$. So \eqref{EqNonDegXDMC} holds true because $\dim X' = \dim X$. Hence we are done.
\end{proof}



\subsection{Technical lemmas}


The following lemma will be useful in the proofs of the desired
bounds. Let for the moment $k$ be an algebraically closed field and
$M\ge 1,n\ge 1$ integers. If $Z$ is a Zariski closed subset of $(\IP_k^n)^M$
we let $\deg Z$ denote the sum of the degrees of all irreducible
components of $Z$ with respect to $\cO(1,\ldots,1)$. 

%\marginpar{\tiny{Noga Alon. Should be a rather elementary result.}}
\begin{lemma}\label{LemmaNAlon} 
  Let $C \subset \IP_k^n$ be an
  irreducible curve defined over $k$ and let $Z \subseteq (\IP_k^n)^M$
  be a Zariski closed subset of $(\IP_k^n)^M$ such that $C^M = C \times
  \cdots \times C \not\subseteq Z$. Then there exists a number $B$,
  depending only on $M, \deg C,$ and $\deg Z$, satisfying the following
  property. If $\Sigma \subset C(k)$ has cardinality $\ge B$, then
  $\Sigma^M = \Sigma \times \cdots \times \Sigma \not\subseteq Z(k)$.
\end{lemma}
\begin{proof}%\Pcomment{I reworded part of this  proof, we don't need flatness or that $C$ is smooth.}
  Let us prove this lemma by induction on $M$. The case $M = 1$ follows
  easily from B\'{e}zout's Theorem.

  Assume the lemma is proved for $1, \ldots, M-1$. Let $q \colon
  (\IP_k^n)^M \rightarrow \IP_k^n$ be the projection to the first
  factor.

  The number of irreducible components of $Z \cap C^M$ and their degrees
  are bounded from above in terms of $M,\deg C,$ and $\deg Z$ by
  B\'ezout's Theorem applied to the Segre embedding.
  Let $Z'$ be the union of all  irreducible
  components $Y$ of $Z \cap C^M$ with $\dim q(Y)\ge 1$, let $Z''$ be the
  union of all other irreducible components.
  % Observe that $q(Y)\subset C$ and $\dim Y \le M-1$ by hypothesis.

  Note that $q(Z')\subset C$. For all $P \in C(k)$ the fiber
  $q|_{Z'}^{-1}(P)=Z'\cap (\{P\}\times (\IP_k^n)^{M-1})$ has
  dimension at most
  $\dim Z'-1\le M-2$. So the projection of $q|_{Z'}^{-1}(P)$ to the final
  factors $(\IP_k^n)^{M-1}$ does not contain $C^{M-1}$. By B\'ezout's
  Theorem the degree of this projection is bounded in terms of $M,\deg
  C,$ and $\deg Z$. We apply the induction hypothesis to the projection
  of $q|_Y^{-1}(P)$ to $(\IP_k^n)^{M-1}$ and obtain a number $B'$,
  depending only on $M,\deg C,$ and $\deg Z$ satisfying the following
  property. If $\Sigma \subseteq C(k)$ has cardinality $\ge B'$, then
  $\{P\} \times \Sigma^{M-1} \not\subseteq Z'(k)$ for all $P \in C(k)$.
  % In particular, $\Sigma^M\not \subset Y(k)$, which concludes this case.

  Now $\dim q(Z'') = 0$, so $q(Z'')$ is a finite set of cardinality
  at most $B''$, the number of irreducible components of
  $Z\cap C^M$. 

  The lemma follows with $B=\max\{B',B''+1\}$.  
\end{proof}

In the next lemma we use the Faltings--Zhang morphism in a single
abelian variety $A$, \textit{i.e.}, $\mathscr{D}_M \colon
A^{M+1}\rightarrow A^M$ defined by $(P_0,\ldots,P_M)\mapsto
(P_1-P_0,\ldots,P_M-P_0)$. 

\begin{lemma}\label{LemmaPacketsAlternative3}
  Let $A$ be an abelian variety defined over $\overline{\IQ}$ and
  suppose $C$ is a smooth curve of genus $g\ge
  2$ contained in $A$. If $Z$ is an irreducible Zariski
  closed and proper subset of $\mathscr{D}_M(C^{M+1})$, 
  then
%in $A^M$, then
  \begin{equation*}
    \#\{ P \in C(\overline{\IQ}) : (C-P)^M \subset Z
%    \subsetneq \mathscr{D}_M(C^{M+1})
    \} \le 84(g-1).
  \end{equation*}
\end{lemma}
\begin{proof}
For simplicity denote by $\Xi = \{ P \in C(\overline{\IQ}) : (C-P)^M
\subset Z \}$. % \subsetneq \mathscr{D}_M(C^{M+1}) \}$.
%
Fix $P_0\in \Xi$. It suffices to prove that there are only $84(g-1)$
possibilities for $P_1 - P_0$ when $P_1$ runs over $\Xi$.

Say $P_1 \in \Xi$ and let $i\in \{0,1\}$. 
Note that
$Z\subsetneq \mathscr{D}_M(C^{M+1})$ and so
 $\dim Z < \dim \mathscr{D}_M(C^{M+1})\le M+1$, 
as $\mathscr{D}_M(C^{M+1})$ irreducible. 
Note that $(C-P_i)^M\subset Z$, and so both have dimension $M$.
As $Z$ is irreducible we find 
 $(C - P_i)^M=Z$ for $i=0$ and $i=1$.

Applying the first projection  $A^M\rightarrow A$ yields $C-P_1=C-P_0$. In other words, $P_1-P_0$
stabilizes $C$. By Hurwitz's Theorem \cite{Hurwitz1892}, a smooth curve over $\IQbar$ of
genus $g\ge 2$ has at most $84(g-1)$ automorphisms.
Hence we are done.
\end{proof}

%%% Local Variables:
%%% TeX-master: "main"
%%% End:
